Trong cơ học cổ điển, bài toán đạn pháo là một mô hình cơ bản mô tả chuyển động của chất điểm bị ném xiên một góc nhất định so với phương ngang, dưới tác dụng của trọng lực. Ở điều kiện lý tưởng không có lực cản không khí, có thể dễ dàng dẫn ra phương trình quỹ đạo một cách giải tích là một hình parabol, từ đó có thể xác định các thông số đặc trưng như tầm xa, độ cao cực đại hay thời gian bay.

Tuy nhiên, phạm vi áp dụng của mô hình lý tưởng thường bị giới hạn trong các ứng dụng thực tế. Sự tương tác giữa chất điểm và môi trường gây ra một lực cản phụ thuộc phi tuyến vào vận tốc, làm thay đổi bản chất của hệ động lực học. Lực cản không chỉ dẫn đến sự sụt giảm đáng kể về tầm xa mà còn phá vỡ tính đối xứng của quỹ đạo, khiến việc tìm nghiệm giải tích vô cùng phức tạp và không còn khả thi vì nó vượt ngoài khả năng biểu diễn của các hàm sơ cấp thông thường.

Để giải quyết tính phi tuyến của hệ, nghiên cứu này tiếp cận bài toán ném xiên thông qua việc xây dựng hệ phương trình vi phân thường bậc nhất và sử dụng phương pháp Runge--Kutta bậc 4 (RK4), một phương pháp xấp xỉ nghiệm của hệ phương trình vi phân thường có độ ổn định và chính xác cao, để mô phỏng quỹ đạo của chất điểm khi chịu tác động của môi trường tồn tại lực cản không khí. Cụ thể, lực cản có dạng tổng quát \(\mathbf{f} = -km |v|^n \frac{\mathbf{v}}{|v|}\) trong đó phần $\frac{\vec{v}}{|v|}$ là vector đơn vị theo hướng vận tốc, đảm bảo lực cản luôn tác dụng ngược chiều chuyển động. Trên có sở đó, chúng tôi tập trung vào ba trường hợp điển hình vận tốc nhỏ $n = 1$, vận tốc trung bình $n = 3/2$ và $n = 2$ cho trường hợp vận tốc cao. Với hệ số lực cản giả định $k = 0.8$, các kết quả thu được sẽ làm rõ sự sai khác về quỹ đạo giữa mô hình lý tưởng và trường hợp có lực cản tại nhiều góc bắn khác nhau, từ đó cho thấy hiệu quả của phương pháp số trong việc phân tích các hệ động lực học phi tuyến.
